\documentclass[11pt]{article}

\usepackage[margin=1in]{geometry}
\usepackage{amsmath}
\usepackage{amssymb}
\usepackage{booktabs}
\usepackage{hyperref}
\usepackage{enumitem}
\usepackage{listings}

\title{gPCD Geometry Dynamics}
\author{}
\date{\today}

\begin{document}
\maketitle
\section{Geo Dynamics Conventions}

\subsection{Purpose}

The Geo path is a clean restart of the collision dynamics code.  Its purpose is
to build a Python implementation whose structure can be transferred to GLSL
with minimal translation.

\subsection{File Boundaries}

\texttt{GeoDynamics.py} is the GLSL-portable dynamics layer.  It should contain
only functions that correspond directly to shader-side dynamics functions.

\texttt{GeoBase.py} is the Python harness.  It may load configuration files,
create Python data structures, build particle arrays, set push constants,
iterate particles, and call the Geo dynamics functions.

\texttt{Reporting.py} is reserved for reporting and diagnostic output.  It
should not own dynamics decisions.

\subsection{Naming Convention}

Functions intended to map to GLSL should use GLSL-like names and signatures
where practical.  Current examples are:

\begin{verbatim}
GeoAddParticleContact(SourceID, TargetID)
GeoProcessCollisions(SourceID)
GeoMoveParticle(SourceID, positionBuffer, dt)
isParticleContact(Frame, SourceID, TargetID, positionBuffer)
\end{verbatim}

\texttt{GeoDynamics.py} should avoid Python-only helper names unless the helper
is also expected to exist in GLSL.

Before adding or renaming any Geo field, buffer, function, or constant, check
the GLSL source files first.  The GLSL names are the naming authority for the
Geo path:

\begin{verbatim}
C:\_DJ\gPCD\vulkan\sim\common\*.glsl
C:\_DJ\gPCD\vulkan\sim\BoundarySpheresMotion\BoundarySpheresMotionP.comp
\end{verbatim}

Python names should stay coordinated with those files whenever possible.  If a
Python-only harness name is needed, it belongs in \texttt{GeoBase.py},
\texttt{GeoRunner.py}, or \texttt{Reporting.py}, not in
\texttt{GeoDynamics.py}.

\subsection{Data Convention}

Particle records use field access:

\begin{verbatim}
Particle.vx
Particle.radius
\end{verbatim}

rather than dictionary access:

\begin{verbatim}
Particle["vx"]
\end{verbatim}

The Python structures should be shaped to resemble GLSL buffers and push
constants.  Python-only configuration parsing must not leak into
\texttt{GeoDynamics.py}.

When adding particle position, velocity, radius, collision-list, push-constant,
or error fields, first look for the corresponding GLSL structure or function
name in the common shader files.  For example, position double buffering should
prefer names coordinated with the GLSL particle buffers and
\texttt{ShaderFlags.positionBuffer}, rather than inventing unrelated Python-only
names.

\subsection{Overlap Area Convention}

Overlap area is a geometric contact property.  For a particle-particle contact,
the source-side and target-side overlap areas are the same geometric area, even
when the two particles have different velocities.  Therefore, once the overlap
area is computed for a source-target pair, the target particle's overlap area is
known from the same calculation.

Velocity may affect momentum transfer, phase, or future position, but it does
not make the instantaneous geometric overlap area different for the two
particles in the same contact.

\subsection{Instantaneous Geometric Goal}

The long-term Geo dynamics target is an instantaneous, purely geometric
collision model.

Instantaneous means the collision response for a frame is computed from the
current frame's particle state only:

\begin{itemize}
    \item current positions;
    \item current velocities;
    \item current radii and masses;
    \item current overlap geometry;
    \item current source-owned contact list for this frame.
\end{itemize}

Purely geometric means the response should be derived from contact geometry,
such as center distance, normal direction, penetration depth, overlap area, and
the distribution of overlap area across a source particle's simultaneous
contacts.  Energy calculations and energy-named dynamics variables should be
reserved for diagnostics only.

The model should not require persistent collision-episode memory such as:

\begin{itemize}
    \item stored contact momentum reservoirs;
    \item immutable first-contact velocity;
    \item persistent compression/rebound phase state;
    \item remembered alpha-zero values;
    \item global wall or contact reservoirs shared by particles.
\end{itemize}

The only persistent particle state should be ordinary particle state, such as
position, velocity, radius, mass, and explicitly approved per-particle fields.
The current frame may build temporary contact information, but that temporary
contact information should be safe to rebuild on the next frame.

\subsection{Instantaneous Frame Snapshot Rule}

The collision response for a frame must not depend on velocity changes induced
earlier in the same frame by other contacts.  Those changes are schedule
artifacts, not instantaneous state.  If two executions see the same frame-start
positions, velocities, radii, masses, contact geometry, and approved
source-owned carried state, they must compute the same dynamics result
independent of source invocation order.

This means collision-induced velocity changes are not a separate dynamics
input.  A frame is treated as one possible instantaneous state that could have
come from many prior histories.  The response may use:

\begin{itemize}[leftmargin=*]
    \item frame-start positions for contact geometry;
    \item frame-start velocities for relative normal velocity and available
    momentum;
    \item current radii, masses, and particle-owned stiffness;
    \item current-frame contact sets and overlap-area weights;
    \item explicitly approved source-owned carried state, such as the
    contact-slot internal normal momentum value.
\end{itemize}

The response must not use:

\begin{itemize}[leftmargin=*]
    \item a target velocity already modified by another source invocation in
    the same frame;
    \item an intra-frame contact response as if it were collision history;
    \item source-order side effects to decide compression, return, or release.
\end{itemize}

Python may satisfy this rule with a frame-start velocity snapshot such as
\texttt{VelRadFrame}.  GLSL must provide an equivalent frame-start read path
before multi-contact behavior can be considered shader-equivalent.

\subsection{Approved Source-Owned Internal Momentum Exception}

The Geo path may carry one source-owned internal normal momentum value per
active contact slot.  This is an explicit exception to the fully stateless
geometric target because the current overlap-area impulse needs a ledger to
avoid releasing more momentum than it removed during compression.

The exception is constrained:

\begin{itemize}[leftmargin=*]
    \item the value is owned by the source particle's contact slot;
    \item it is not a shared pair reservoir;
    \item it must not write the target particle;
    \item it may be used only to cap source-owned rebound/release impulse;
    \item it must remain representable in the GLSL \texttt{ContactState}
    structure.
\end{itemize}

In the current GLSL-shaped layout, the carried value is:

\begin{verbatim}
P[SourceID].contacts[slot].aux.z
\end{verbatim}

This field stores source-owned internal normal momentum.  It exists to prevent
symmetric velocity drift while preserving the source-owned shader execution
model.

\subsection{Model-Preservation Warning}

Before editing code, any proposed change that would move the Geo path away from
the instantaneous, purely geometric target must be called out explicitly.

Examples of changes that require warning before implementation:

\begin{itemize}
    \item adding persistent per-contact memory for dynamics;
    \item adding a reservoir whose value survives after the frame;
    \item using first-contact velocity as a long-lived dynamics input;
    \item using phase history to decide future dynamics;
    \item adding global writable memory shared by all particles;
    \item adding a wall-specific dynamics path instead of expressing the wall as
    contact geometry.
\end{itemize}

This warning does not forbid such changes forever, but it makes the model cost
visible before implementation.

\subsection{Documenting Problem Evolution}

When a dynamics rule changes, the reason for the change should be recorded in
the problem-evolution document.  That document is the design history for the
Geo path.  It should explain how the problem was observed, which fixes were
rejected, which rule or exception was accepted, and what result verified the
change.

New entries should follow this pattern:

\begin{itemize}[leftmargin=*]
    \item identify the observed failure by test or symptom;
    \item include the commit hash associated with the change;
    \item describe rejected fixes when they would violate the GLSL execution
    model;
    \item state any convention, rule, or exception that was added;
    \item describe the implemented algorithm in equations or direct procedural
    terms;
    \item record the verification result that shows the issue was resolved.
\end{itemize}

The purpose is not only to document the current code.  It is also to preserve
the assumptions behind that code so future work can deliberately keep, revise,
or remove those assumptions.

\subsection{Control Flow Shape}

The Python control-flow harness should match the compute shader shape:

\begin{verbatim}
for each SourceID:
    build source-owned contact list
    GeoProcessCollisions(SourceID)
    GeoMoveParticle(SourceID, positionBuffer, dt)
\end{verbatim}

The current Python harness may use a direct all-pairs scan while the shader uses
cell/corner occupancy lookup.  That difference belongs in \texttt{GeoBase.py},
not \texttt{GeoDynamics.py}.

\subsection{Sequential Stage And Narrow Function Rule}

Geo frame processing should use a sequential, data-oriented stage chain.  The
stage chain owns the visible decisions about what kind of work happens next.
The narrow functions called by the chain should each perform one kind of
activity.

For example, a frame-level chain may decide:

\begin{verbatim}
build frame snapshot
build particle contact lists
build wall contact lists
resolve particle contacts
resolve wall contacts
move particles
report diagnostics
\end{verbatim}

The functions inside those stages should stay narrow:

\begin{itemize}[leftmargin=*]
    \item \texttt{isParticleContact(...)} answers only whether two particles
    are touching in the current snapshot;
    \item \texttt{isWallContact(...)} answers only whether a particle touches a
    wall or wall particle in the current snapshot;
    \item \texttt{GeoAddParticleContact(...)} records a particle contact only;
    \item \texttt{GeoAddWallContact(...)} records a wall contact only;
    \item geometry functions compute geometry only;
    \item impulse functions compute impulse values only;
    \item apply functions write the source-owned state only.
\end{itemize}

The narrower the split, the more situations can share the same downstream
dynamics.  Particle-particle contacts, wall contacts, and wall ghost-particle
contacts should differ only where they truly differ.  Usually that is contact
determination and contact creation:

\begin{verbatim}
isParticleContact(...)
isWallContact(...)
GeoAddParticleContact(...)
GeoAddWallContact(...)
\end{verbatim}

After a contact has been expressed as source-owned contact state, later stages
should be shared whenever possible:

\begin{verbatim}
GeoContactGeometry(...)
GeoContactContext(...)
GeoContactWeight(...)
GeoCompressionImpulse(...)
GeoReturnImpulse(...)
GeoApplySourceImpulse(...)
\end{verbatim}

This prevents the code from growing separate particle and wall dynamics paths
when both contacts can use the same response rule.  If a situation needs a
different narrow function, add that narrow function at the point of real
difference instead of cloning a broad stage.

Do not add an umbrella stage when it only hides useful frame flow.  For
example, this is less clear:

\begin{verbatim}
GeoBuildContactLists()
\end{verbatim}

than leaving the meaningful decisions visible in the main chain:

\begin{verbatim}
GeoBuildParticleContactLists()
GeoBuildWallContactLists()
\end{verbatim}

Umbrella functions are allowed only when they remove real duplication or hide
mechanics that are not relevant to understanding the frame sequence.

This style makes the program easier to step through, easier to debug, and
easier to learn.  If a contact list is wrong, the defect should be local to the
contact-detection or contact-recording stage.  If an impulse is wrong, the
defect should be local to the context, weighting, or impulse stage.  Decisions
about which narrow functions are called should remain visible in the stage
chain instead of being hidden inside broad helper functions.

This rule applies to the Geo dynamics path and its frame harness.  Unrelated
Python utilities do not need to be rewritten to match this style.

\subsection{No Pairwise Write-Back Shortcut}

The Geo Python implementation must not use a pairwise dynamics shortcut that
cannot be exported to the current GLSL execution model.

The active shader shape is source-owned: one invocation owns one source
particle.  Therefore a Geo dynamics function may write only the current source
particle's state.  It must not solve a pair by directly updating both
particles:

\begin{verbatim}
solve_pair(source, target):
    P[source].velocity = ...
    P[target].velocity = ...   // forbidden in GeoDynamics.py
\end{verbatim}

This rule applies even when a pairwise update would be physically cleaner in
Python.  Python is not allowed to become a separate reference model that fixes
collisions in a way Vulkan cannot reproduce.  If the shader cannot safely do
the write pattern, the Python Geo path cannot do it either.

Allowed alternatives must preserve source-owned execution.  Examples include:

\begin{itemize}[leftmargin=*]
    \item computing a source-owned response from source-owned contact data;
    \item using read-only frame-start snapshots of target state;
    \item accumulating source-owned diagnostic values;
    \item designing an explicit GLSL-compatible staged pass before adding the
    equivalent Python implementation.
\end{itemize}

Any proposal that needs one contact solve to write both particles must first
change the GLSL architecture, not just the Python model.

\subsection{Dynamics Versus Harness}

\texttt{GeoDynamics.py} should not:

\begin{itemize}
    \item read cfg files;
    \item create particle lists from cfg files;
    \item know about UI flags;
    \item know about reports or capture files;
    \item contain test harness iteration logic.
\end{itemize}

\texttt{GeoBase.py} may do those things because it is not intended to be a
one-to-one GLSL export unit.

\subsection{Current Standing Rules}

\begin{itemize}[leftmargin=*]
    \item Keep \texttt{GeoDynamics.py} one-to-one with GLSL export.
    \item Add contact handling only where it matches the shader model.
    \item Add reporting only through \texttt{Reporting.py}, not through
    dynamics.
    \item Record dynamics-rule changes in \texttt{ProblemEvolution.tex}.
\end{itemize}

\section{Particle Geometry} \label{New Fundamental}

The Geo model treats collision response as a current-frame geometric operation.
The response is computed from the state that exists in the current simulation
step:

\begin{itemize}[leftmargin=*]
    \item current particle positions,
    \item current particle velocities,
    \item current contact normal,
    \item current penetration depth,
    \item current overlap area,
    \item and the source particle's simultaneous contact areas.
\end{itemize}

The model should not require a remembered first-contact velocity, a stored
turnaround depth, or a persistent compression/rebound phase.  Those quantities
may be useful diagnostics, but they are not required dynamics state in the
instantaneous geometric model.

\section{Current Relative Normal Velocity}

For a source particle $s$ and target particle $t$, let the current contact
normal be:

\begin{equation}
\begin{aligned}
\hat{n}_{s,t}
=
\frac{x_t - x_s}{\lVert x_t - x_s \rVert}
\end{aligned}
\end{equation}

The current relative normal velocity is:

\begin{equation}
\begin{aligned}
v_n
=
(v_t - v_s) \cdot \hat{n}_{s,t}
\end{aligned}
\end{equation}

With this sign convention:

\begin{itemize}[leftmargin=*]
    \item $v_n < 0$ means the source and target are closing,
    \item $v_n = 0$ means there is no relative normal motion,
    \item $v_n > 0$ means the source and target are separating.
\end{itemize}

The collision response uses this current-frame value.  It is not the velocity
from the first frame of a collision episode.

\section{Overlap Geometry}

For two circular particles, the current overlap area is:

\begin{equation}
\begin{aligned}
A_{s,t} = A(r_s, r_t, d)
\end{aligned}
\end{equation}

where $r_s$ and $r_t$ are the particle radii and $d$ is the current center
distance.

The overlap area is a geometric property of the contact.  The source-side and
target-side overlap areas are the same area:

\begin{equation}
\begin{aligned}
A_{s,t} = A_{t,s}
\end{aligned}
\end{equation}

Velocities can change the response, but they do not change the instantaneous
geometric overlap area.

\section{Geometric Response}

The current-frame geometric response should be a function of overlap area,
current relative normal velocity, stiffness, and source-area weighting:

\begin{equation}
\begin{aligned}
\Delta v_{s,t}
=
G(A_{s,t}, v_n, Q, w_{s,t})
\end{aligned}
\end{equation}

where:

\begin{itemize}[leftmargin=*]
    \item $A_{s,t}$ is the current overlap area,
    \item $v_n$ is the current relative normal velocity,
    \item $Q$ is the stiffness parameter,
    \item $w_{s,t}$ is the source-owned area weight for this contact.
\end{itemize}

The exact form of $G$ is the next dynamics decision.  The important convention
is that $G$ must be computed from the current frame.  It should not depend on
stored first-contact values.

\section{Source-Owned Area Weighting}

A source particle may contact multiple targets in the same frame.  The source
particle's total current overlap area is:

\begin{equation}
\begin{aligned}
A_s^{total}
=
\sum_{t \in C_s} A_{s,t}
\end{aligned}
\end{equation}

where $C_s$ is the current contact set for source particle $s$.

The area weight for one contact is:

\begin{equation}
\begin{aligned}
w_{s,t}
=
\frac{A_{s,t}}{A_s^{total}}
\end{aligned}
\end{equation}

when $A_s^{total} > 0$.

These weights keep the source particle from responding to each contact as if it
were the only contact.  In a multi-contact frame, the source response must be
distributed across the current overlap areas.

\section{Legacy Turnaround Formulas}

Previous notes used an incoming collision velocity $v_{in}$ and a stored
turnaround depth:

\begin{equation}
\begin{aligned}
\alpha_0
=
\left(\frac{3\mu v_{in}^2}{2Q}\right)^{1/3}
\end{aligned}
\end{equation}

This formula remains useful as a diagnostic or comparison curve.  In the new
Geo model, however, $\alpha_0$ should not be stored as persistent dynamics
state.  If a zero-velocity or turnaround estimate is needed, it should be
derived from the current frame and treated as temporary diagnostic information.

\section{Starting In Contact}

When particles start in contact, the response should still be computed from the
current frame:

\begin{itemize}[leftmargin=*]
    \item current overlap area,
    \item current relative normal velocity,
    \item stiffness,
    \item and current source-area weights.
\end{itemize}

The simulation should not need to infer an earlier collision history.  If the
starting overlap is too deep for the selected geometric response rule, that is a
guardrail or configuration issue, not a reason to create hidden collision
history.

% !TeX root = GeoDynamics.tex
\section{Multi-Particle Contacts}

Our model states that the dynamics of a particle depends only on the area of overlap into another particle. 


Momentum is summed The momentum in collision must equal the amount total momentum per area in rebound. But in multi-contact, a source particle cannot independently apply the full area rule to every contact as if each contact owns the whole source response capacity. That double-counts available momentum. For a particle with multiple active contacts, the source has one current momentum state and multiple overlap areas. The contacts must share the source response by weight.  This is the fundamental source of conservation of the whole system. Therefore, schedule wise, \textbf{source owned contact resolution must work}. In multiple particle contacts the weighing of area, changes the momentum per area.

In the geometric compression model, the momentum at a given frame is computed from the current overlap area, not by integrating every prior step.So it is instantaious. In the geometric compression model, the momentum at a given frame is computed from the current overlap area, not by integrating every prior step.
 

But, 

Currently multi-particle contacts averages the velocity of each particle pair.    
This is not a geometric solution but a velocity prediction. 

For a contact, the change in normal velocity is proportional to the penetration area. 
The momentum response associated with that contact is proportional to mass times that velocity change.
The relative contribution of each contact is proportional to its overlap area divided by the total active overlap area for the source particle.

The total overlap area is the sum of all target particle
overlap areas. From this the total momentum imposed on the target particle can be found.
Taking the ratio of total overlap area to the individual overlap areas gives the fraction
each target particle contributes to the collision.

From section \ref{Fundamental}, we have three relationships:
\begin{enumerate}
\item Velocity/area relation
\item Velocity change proportional to penetration area
\item Force/penetration law
\end{enumerate}

\subsection{Velocity}

The clearest rule for Mutli-Particle Contact is from the second fundamental relation:
\begin{equation}
\begin{aligned}
(v_{in} - v_p) = P A_P
\end{aligned}
\end{equation}

so that :
\begin{equation}
\begin{aligned}
\Delta v = P A_P
\end{aligned}
\end{equation}

That gives the clean area relationship:

\textbf{Velocity change is proportional to penetration/overlap area}

\subsection{Momentum}

For momentum:

\begin{equation}
\begin{aligned}
\Delta p = m \Delta v
\end{aligned}
\end{equation}

so that:


\begin{equation}
\begin{aligned}
\Delta p = m P A_P
\end{aligned}
\end{equation}

So the momentum relationship states:

\textbf{Contact momentum contribution is proportional to overlap area}

For several contacts on the same source particle:

\begin{equation}
\begin{aligned}
A_T = \sum_k A_k
\end{aligned}
\end{equation}

Where each contact’s geometric weight is:
\begin{equation}
\begin{aligned}
w_k = \frac{A_k}{A_T}
\end{aligned}
\end{equation}

Which states that:

\textbf{Each target contact contributes a fraction of the source particle's total contact response according to its overlap-area weight.}

From the local relation:
\begin{equation}
\begin{aligned}
\label{none}
\Delta v = P A_P
\end{aligned}
\end{equation}

and since:
\begin{equation}
\begin{aligned}
\label{none}
\Delta p = m \Delta v
\end{aligned}
\end{equation}

we get:


\begin{equation}
\begin{aligned}
\label{none}
\Delta p = m P A_P
\end{aligned}
\end{equation}

So the simplest area-to-momentum relationship is:

\begin{equation}
\begin{aligned}
\label{none}
\Delta p \propto m A_P
\end{aligned}
\end{equation}

But for contact dynamics, the better mass term is probably not just source mass. For a pair contact, the normal exchange depends on reduced mass:

\begin{equation}
\begin{aligned}
\label{none}
\mu_{ij} = \frac{m_i m_j}{m_i + m_j}
\end{aligned}
\end{equation}

So for a contact $k$ between source particle $i$ and target particle $j$:

\begin{equation}
\begin{aligned}
\label{none}
\Delta p_k = C \mu_{ij} A_k
\end{aligned}
\end{equation}

where:

C = contact stiffness / momentum-per-area proportionality
$A_k$ = overlap area for that contact
Then for all contacts owned by source particle $i$:

\begin{equation}
\begin{aligned}
\label{none}
\Delta p_{total,i}= C \sum_k \mu_{ik} A_{ik}
\end{aligned}
\end{equation}

That gives a geometric total contact momentum scale.

Then each contact gets:

\begin{equation}
\begin{aligned}
\label{none}
w_k =
\frac{\mu_{ik} A_{ik}}
{\sum_j \mu_{ij} A_{ij}}
\end{aligned}
\end{equation}

and:

\begin{equation}
\begin{aligned}
\label{none}
\Delta p_k = w_k \Delta p_{total,i}
\end{aligned}
\end{equation}

But because $w_k$ is built from the same quantity, this simplifies to:


\begin{equation}
\begin{aligned}
\label{none}
\Delta p_k = C \mu_{ik} A_{ik}
\end{aligned}
\end{equation}

The total is just the sum:

\begin{equation}
\begin{aligned}
\label{none}
\Delta p_{total,i}=\sum_k \Delta p_k
\end{aligned}
\end{equation}

So the real decision is the coefficient $C$.

In the old code, the closest thing is:
\begin{verbatim}
overlap_contact_momentum =
    momentum_per_area * overlap_area * inverse_square_weight(center_distance)
\end{verbatim}
So the current practical relationship is:

\begin{equation}
\begin{aligned}
\label{none}
\Delta p_k=P A_k W(d_k)
\end{aligned}
\end{equation}

where:

\begin{equation}
\begin{aligned}
\label{none}
W(d_k) = \frac{1}{\max(d_k^2, \epsilon^2)}
\end{aligned}
\end{equation}

But if we want to move more geometric and less heuristic, the cleaner relationship is:

\begin{equation}
\begin{aligned}
\label{none}
\Delta p_k = P A_k
\end{aligned}
\end{equation}

or with reduced mass:

\begin{equation}
\begin{aligned}
\label{none}
\Delta p_k = P \mu_{ik} A_k
\end{aligned}
\end{equation}

For equal masses, reduced mass is constant, so the weighting reduces to area weighting.

The important correction for multi-particle contact is:

text


Do not average pair velocity predictions.
Compute a contact momentum vector from area for each active contact,
sum those contact momentum vectors,
then apply the resulting source momentum change once.
Vector form:


\begin{equation}
\begin{aligned}
\label{none}
\Delta \mathbf{p}_{i}=-\sum_k\left(P \mu_{ik} A_{ik}\right)\mathbf{n}_{ik}
\end{aligned}
\end{equation}

where $\mathbf{n}_{ik}$ points from source to target.

Then:

\begin{equation}
\begin{aligned}
\label{none}
\Delta \mathbf{v}_{i}=\frac{\Delta \mathbf{p}_{i}}{m_i}
\end{aligned}
\end{equation}

That gives us the missing $\Delta \mathbf{p}_{total}$

\begin{equation}
\begin{aligned}
\label{none}
\Delta \mathbf{p}_{total,i}=-\sum_k P \mu_{ik} A_{ik} \mathbf{n}_{ik}
\end{aligned}
\end{equation}


This is source-owned, geometric, and GLSL-friendly. It is not automatically globally momentum-conserving unless every source computes its matching response consistently, but it removes the current bad average and gives a concrete area-to-momentum rule.







\subsection{Force}

\begin{equation}
\begin{aligned}
\label{frclaw}
F(\alpha) = Q\alpha^2
\end{aligned}
\end{equation}

Equation \ref{frclaw}, is used to predict the zero-velocity depth: 
\begin{equation}
\begin{aligned}
\label{zerodepth}
\alpha_0 =
\left(\frac{3\mu v_{in}^2}{2Q}\right)^{1/3}
\end{aligned}
\end{equation}


\subsection{Multicontact rule}

The multi-particle rule should is stated as:

The change in velocity is proportional to the total collision overlap area. 
\begin{equation}
\begin{aligned}
\label{mc}
\Delta v_k \propto A_k
\end{aligned}
\end{equation}

The contributing weight of each collision partner is the ratio of target overlap to total overlap.

\begin{equation}
\begin{aligned}
\label{mc}
w_k = \frac{A_k}{\sum_j A_j}
\end{aligned}
\end{equation}

The change in momentum is equal to the total momentum times the weight.

\begin{equation}
\begin{aligned}
\label{mca}
\Delta \mathbf{p}_k = w_k \Delta \mathbf{p}_{total}
\end{aligned}
\end{equation}


% !TeX root = GeoDynamics.tex
\section{Contact Prediction Processing}

This section documents the contact prediction process. Contact prediction is
the operation that takes one active contact and produces a source-particle
velocity prediction for that contact. The schedule-wise resolution algorithm
then combines these predictions.

Prediction is source-owned. A prediction may read the target particle and the
source-owned contact state, but it returns only a velocity for the source
particle.

\subsection{Particle Contact Prediction Algorithm}

For a source particle \(i\) and target particle \(j\), the particle contact
prediction algorithm is:

\begin{enumerate}
    \item Read the current contact geometry:
    \[
        \mathbf{n}_{ij},\quad A,\quad d .
    \]
    Here \(\mathbf{n}_{ij}\) points from source particle \(i\) to target
    particle \(j\), \(A\) is the current overlap area, and \(d\) is the current
    center distance.

    \item Compute the current relative normal velocity:
    \[
        u_{ij}
        =
        \left(
            \mathbf{v}_j
            -
            \mathbf{v}_i
        \right)
        \cdot
        \mathbf{n}_{ij}.
    \]

    \item Choose the provisional phase:
    \[
        u_{ij} < 0 \Rightarrow \text{compression},
        \qquad
        u_{ij} \ge 0 \Rightarrow \text{rebound}.
    \]

    \item Read the source-owned contact state. If the state contains stored
    zero-velocity geometry, update the phase using the current overlap and
    current relative normal velocity. If there is no zero-velocity geometry and
    the state is not already rebound, the prediction is invalid.

    \item Read the first-contact source and target velocities from the contact
    state:
    \[
        \mathbf{v}_{i}^{0},
        \qquad
        \mathbf{v}_{j}^{0}.
    \]
    If these are missing, the current velocities are used as fallback values.

    \item Construct a two-particle Neo prediction using the current geometry,
    first-contact velocities, masses, radii, and stored zero-velocity geometry.

    \item Compute the Neo report at the current overlap. The report contains
    both compression and rebound velocities:
    \[
        \mathbf{v}_{i,compression}^{\ast},
        \mathbf{v}_{j,compression}^{\ast},
        \mathbf{v}_{i,rebound}^{\ast},
        \mathbf{v}_{j,rebound}^{\ast}.
    \]

    \item Update the diagnostic internal momentum values stored in the contact
    state.

    \item Select the source and target predicted velocities from the report:
    \[
        \text{compression}
        \Rightarrow
        \left(
            \mathbf{v}_{i,compression}^{\ast},
            \mathbf{v}_{j,compression}^{\ast}
        \right),
    \]
    \[
        \text{rebound}
        \Rightarrow
        \left(
            \mathbf{v}_{i,rebound}^{\ast},
            \mathbf{v}_{j,rebound}^{\ast}
        \right).
    \]

    \item If the phase is rebound, apply the outward escape rule to the pair of
    predicted velocities. This enforces a small minimum outward relative normal
    velocity:
    \[
        u_{out} \ge u_{min}.
    \]

    \item If the phase is rebound, apply the internal-contact momentum feedback
    term to the source prediction.

    \item Preserve the source particle's current tangential velocity. Only the
    predicted normal component is transferred into the returned source
    velocity:
    \[
        \mathbf{v}_{i}^{return}
        =
        \mathbf{v}_i
        -
        \left(
            \mathbf{v}_i \cdot \mathbf{n}_{ij}
        \right)
        \mathbf{n}_{ij}
        +
        \left(
            \mathbf{v}_{i}^{\ast} \cdot \mathbf{n}_{ij}
        \right)
        \mathbf{n}_{ij}.
    \]

    \item Return:
    \[
        \mathbf{v}_{i}^{0},
        \qquad
        \mathbf{v}_{i}^{return}.
    \]
\end{enumerate}

\subsection{Wall Contact Prediction Algorithm}

Wall contact prediction uses the same contact prediction machinery with a wall
ghost particle. The wall ghost is fixed and is not integrated.

For source particle \(i\) and wall \(w\), the wall prediction algorithm is:

\begin{enumerate}
    \item Build the wall ghost contact geometry:
    \[
        \mathbf{n}_{w},\quad A,\quad d.
    \]
    The normal \(\mathbf{n}_{w}\) points from the particle toward the wall
    ghost.

    \item Choose the velocity accumulator to process. If no accumulator is
    provided, use the particle's current velocity:
    \[
        \mathbf{v}_{acc} = \mathbf{v}_i .
    \]

    \item Compute the wall normal velocity:
    \[
        u_w
        =
        \mathbf{v}_{acc}
        \cdot
        \mathbf{n}_{w}.
    \]

    \item Choose the provisional wall phase:
    \[
        u_w > 0 \Rightarrow \text{compression},
        \qquad
        u_w \le 0 \Rightarrow \text{rebound}.
    \]

    \item Read the wall contact state and update the phase against the stored
    zero-velocity geometry. The wall update uses:
    \[
        u_{relative} = -2u_w
    \]
    so the wall ghost path follows the same pair-style phase logic.

    \item Call the generic contact prediction operation with the real particle
    as source and the wall ghost as target.

    \item If the phase is rebound and zero-velocity geometry exists, replace
    the generic pair rebound prediction with the wall-specific rebound
    prediction. The wall-specific prediction reflects the source particle
    against the fixed wall:
    \[
        \mathbf{v}_{full}
        =
        \mathbf{v}^{0}
        -
        2
        \left(
            \mathbf{v}^{0}\cdot\mathbf{n}_{w}
        \right)
        \mathbf{n}_{w}.
    \]

    \item Preserve the current tangential component of the accumulator and
    replace only its normal component:
    \[
        \mathbf{v}_{return}
        =
        \mathbf{v}_{acc}
        -
        \left(
            \mathbf{v}_{acc}\cdot\mathbf{n}_{w}
        \right)
        \mathbf{n}_{w}
        +
        \left(
            \mathbf{v}^{\ast}\cdot\mathbf{n}_{w}
        \right)
        \mathbf{n}_{w}.
    \]

    \item Return:
    \[
        \mathbf{v}_{acc},
        \qquad
        \mathbf{v}_{return}.
    \]
\end{enumerate}

\subsection{Phase Update Algorithm}

The prediction process depends on whether the contact is compressing or
rebounding. The phase update algorithm is:

\begin{enumerate}
    \item Classify the current motion mode from relative normal velocity:
    \[
        u < -\epsilon \Rightarrow \text{accumulating},
    \]
    \[
        u > \epsilon \Rightarrow \text{releasing},
    \]
    \[
        |u| \le \epsilon \Rightarrow \text{blocked}.
    \]

    \item Compare the current overlap area \(A\) to the stored zero-velocity
    overlap area \(A_0\). With tolerance \(q\):
    \[
        A \ge A_0(1-q)
        \Rightarrow
        \text{zero reached}.
    \]

    \item If the contact reversed outward before reaching \(A_0\), promote the
    zero-velocity state to the current geometry:
    \[
        A_0 \leftarrow A,
        \qquad
        d_0 \leftarrow d.
    \]

    \item If zero has been reached, set the phase to rebound.

    \item Otherwise, if motion is accumulating, set the phase to compression.

    \item Otherwise, if motion is releasing or blocked, keep or set the phase
    to rebound.
\end{enumerate}

\subsection{Pseudocode}

\begin{verbatim}
particle_contact_prediction(source i, target j):
    contact = particle_contact(i, j)
    if contact does not exist:
        return invalid

    normal, overlap_area, center_distance = contact
    relative_normal_velocity = dot(v[j] - v[i], normal)

    phase =
        compression if relative_normal_velocity < 0
        rebound otherwise

    state = source_owned_contact_state(i, j)
    zero = zero_velocity_geometry(state)

    if zero exists:
        update_phase(state, overlap_area, center_distance,
                     relative_normal_velocity)
        phase = state.phase
    else if state.phase is not rebound:
        return invalid

    prediction =
        generic_contact_prediction(i, j, contact, state, phase)

    return prediction


generic_contact_prediction(source i, target j, contact, state, phase):
    read first-contact source and target velocities
    build temporary two-particle Neo model
    compute current-overlap Neo report
    update diagnostic internal momentum state

    if phase is compression:
        source_velocity = report.source_compression_velocity
        target_velocity = report.target_compression_velocity
    else:
        source_velocity = report.source_rebound_velocity
        target_velocity = report.target_rebound_velocity
        source_velocity, target_velocity =
            apply_outward_escape(source_velocity, target_velocity)
        source_velocity =
            apply_internal_contact_momentum_feedback(source_velocity)

    source_velocity =
        preserve_current_tangent_and_use_predicted_normal(source_velocity)

    return first_contact_source_velocity, source_velocity


wall_contact_prediction(source i, wall w, accumulator velocity):
    contact = wall_ghost_contact(i, w)
    if contact does not exist:
        return invalid

    normal_velocity = dot(accumulator_velocity, wall_normal)
    phase =
        compression if normal_velocity > 0
        rebound otherwise

    update wall phase using relative velocity -2 * normal_velocity

    prediction =
        generic_contact_prediction(i, wall_ghost, contact, state, phase)

    if phase is rebound:
        prediction = wall_specific_rebound_prediction(...)

    prediction =
        preserve_accumulator_tangent_and_use_predicted_normal(prediction)

    return accumulator_velocity, prediction
\end{verbatim}

\subsection{Known Questions}

The prediction process is source-owned, but it still contains several choices
that affect dense-contact behavior:
\begin{itemize}
    \item when zero-velocity geometry should be promoted to the current overlap,
    \item whether outward escape should be applied inside each contact
          prediction or only after contact predictions are combined,
    \item whether internal-contact momentum feedback should be applied before
          or after schedule-wise combination,
    \item whether wall prediction should use pair-style ghost prediction,
          wall-specific prediction, or a single unified expression.
\end{itemize}

% !TeX root = GeoDynamics.tex
\section{Schedule-Wise Contact Resolution}

This section documents the current schedule-wise contact combination rule. The
goal is to make the Python path match the GPU execution model as closely as
possible before porting the dynamics to GLSL.

The rule is source-owned:
\begin{itemize}
    \item one source particle computes its own velocity update,
    \item the source particle may read target particles and wall contact state,
    \item the source particle does not write target particle velocity,
    \item pair contacts, wall contacts, and final escape correction are applied
          in a fixed schedule.
\end{itemize}

This replaces the older cluster solve for the GLSL-portable path. The cluster
solver is retained only as a diagnostic comparator.

\subsection{Definitions}

For a source particle \(i\), let its current velocity be
\[
    \mathbf{v}_i .
\]

For each active particle contact \(c\), the contact resolver updates the
contact state and returns a source-owned contact momentum vector:
\[
    \Delta \mathbf{p}_{i,c}.
\]
This vector is computed from the contact's geometric/internal momentum scalar
and the contact normal.  The scheduler sums these momentum vectors and applies
one source velocity update.

\subsection{Algorithm}

For each source particle \(i\), the current schedule-wise algorithm is:

\begin{enumerate}
    \item Start with the particle's current velocity:
    \[
        \mathbf{v}_{base} = \mathbf{v}_i .
    \]

    \item Initialize the pair-contact momentum accumulator:
    \[
        \Delta \mathbf{p}_{sum} = (0,0).
    \]

    \item For each active particle-particle contact owned by source particle
    \(i\):
    \begin{enumerate}
        \item Update the source-owned contact state.

        \item If this is the first active pair contact, set:
        \[
            \mathbf{v}_{base} = \mathbf{v}_{i,c}^{0}
        \]
        where \(\mathbf{v}_{i,c}^{0}\) is the source velocity at first contact.

        \item Compute the contact momentum vector:
        \[
            \Delta \mathbf{p}_{i,c}
            =
            -p_{i,c}\mathbf{n}_{i,c}.
        \]

        \item Accumulate the contact momentum vector:
        \[
            \Delta \mathbf{p}_{sum}
            \leftarrow
            \Delta \mathbf{p}_{sum}
            +
            \Delta \mathbf{p}_{i,c}.
        \]
    \end{enumerate}

    \item If there are active pair contacts, apply the summed momentum once:
    \[
        \mathbf{v}_{pair}
        =
        \mathbf{v}_{base}
        +
        \frac{\Delta \mathbf{p}_{sum}}{m_i}.
    \]
    If there are no pair contacts, use:
    \[
        \mathbf{v}_{pair} = \mathbf{v}_i .
    \]

    \item Apply wall contacts sequentially to \(\mathbf{v}_{pair}\). For each
    active wall contact \(w\), compute the wall-predicted normal velocity and
    replace only the normal component of the current accumulator:
    \[
        \mathbf{v}_{wall}
        =
        \mathbf{v}_{acc}
        +
        \left(
            v_{n,w}^{\ast}
            -
            \mathbf{v}_{acc}\cdot\mathbf{n}_w
        \right)
        \mathbf{n}_w .
    \]
    After each wall contact:
    \[
        \mathbf{v}_{acc} \leftarrow \mathbf{v}_{wall}.
    \]

    \item Apply the pair rebound escape scheduler. For each active rebound pair
    contact, compute the desired outward relative normal velocity. If the
    current relative normal velocity is too small, adjust only the source
    particle accumulator:
    \[
        \mathbf{v}_{acc}
        \leftarrow
        \mathbf{v}_{acc}
        -
        \frac{1}{2}
        \left(
            u_{desired}
            -
            u_{current}
        \right)
        \mathbf{n}_{ij}.
    \]

    \item Write the final accumulator back to the source particle:
    \[
        \mathbf{v}_i \leftarrow \mathbf{v}_{acc}.
    \]
\end{enumerate}

\subsection{Pseudocode}

\begin{verbatim}
for each source particle i:
    base_velocity = current_velocity(i)
    accumulated_momentum = (0, 0)
    has_pair = false

    for each active particle contact c owned by i:
        contribution = pair_momentum_contribution(i, c)
        if contribution is invalid:
            continue

        if not has_pair:
            base_velocity = contribution.base_velocity
            has_pair = true

        accumulated_momentum += contribution.contact_momentum_vector

    if has_pair:
        pair_velocity =
            base_velocity + accumulated_momentum / mass(i)
    else:
        pair_velocity = current_velocity(i)

    final_velocity = pair_velocity

    for each active wall contact w owned by i:
        final_velocity = wall_prediction(i, w, final_velocity)

    for each active rebound pair contact c owned by i:
        final_velocity =
            apply_pair_escape_if_needed(i, c, final_velocity)

    velocity(i) = final_velocity
\end{verbatim}

\subsection{Known Question}

The current combination rule sums source-owned contact momentum vectors.  This
matches the geometric multi-contact convention more closely than averaging
velocity predictions, but it is still not a coupled global solve.  Dense contact
groups may still require limiting, weighting, or another source-owned
geometric rule.

\input{porting.tex}
% !TeX root = GeoDynamics.tex
\section{Drift calculations}

Drift calculation is an accumulation of momentum changes during a simulation. The drift is:


\begin{equation}
\begin{aligned}
\mathbf{p}_{CurrentTotalMomentum} = \mathbf{p}_{ParticleVelocityMomentum} \\ + \mathbf{p}_{ ParticleSideInternalMomentum} + \mathbf{p}_{WallGhostInternalMomentum} + \mathbf{p}_{WallVelocityMomentumReservoir}
\end{aligned}
\end{equation}

Where velocity momentum, $\mathbf{p}_{ParticleVelocityMomentum} $, is the sum of all the particles translational momentum. Internal momentum, $\mathbf{p}_{ParticleSideInternalMomentum}$. Wall momentum \\ $\mathbf{p}_{ WallGhostInternalMomentum}$ is the value of the wall internal momentum. And \\ $\mathbf{p}_{WallVelocityMomentumReservoir}$ is the pseudo wall reaction velocity.

The total starting momentum $\mathbf{p}_{InitialTotalMomentum}$, minus the current total momentum \\  $\mathbf{p}_{CurrentTotalMomentum}$ is drift $\mathbf{p}_{Drift}$.

Reporting variable perform a running sum where as each particles momentum values are recorded by adding them to previous reporting values. This means for each particle, at a time step, the kinetic momentum is added to the global variable $\mathbf{p}_{ ParticleVelocity Momentum}$, and likewise for the other variables. They are independently assigned to reporting variables. At the end of the time step the filled reporting variables are used to calculate the drift.  

Drift tells us how much the total system is varying from the current system and shows us the accumulation/system errors. The  accumulation/system errors become very important in large scale simulation of millions of particles.

The model of wall internal momentum may not match the actual dynamics but it is viewed as follows: The wall/ghost internal momentum ,$\mathbf{p}_{WallGhostInternalMomentum}$ is equal and opposite to the particle-side,$\mathbf{p}_{ParticleSideInternalMomentum}$  wall-contact internal momentum. 

The model of particle internal momentum may not match the actual dynamics but it is viewed as follows:Each particles internal momentum accumulates during compression and feeds back during rebound. During compression it is a vector opposite the velocity vector. In this way the internal momentum is fed back in the same direction as rebound. 

Particle-particle collisions result in four offsetting vectors, two opposing velocity momentum vectors and two opposing internal momentum vectors. Particle-wall collisions differ mathematically since the wall does not move. There are then only three vectors - two internal momentum vectors and only one velocity momentum vector. This leaves an imbalance of one velocity momentum vector. To account for this we create a wall-side velocity momentum reservoir $\mathbf{p}_{WallVelocityMomentumReservoir}$.





\section{Minimal GLSL contact state}

In Python it is useful to store extra diagnostic values.  In GLSL the goal is
to store only the persistent values that cannot be safely reconstructed from the
current frame.

The first implementation assumes:

\begin{itemize}[leftmargin=*]
    \item natural contacts begin outside collision,
    \item contacts that exist at frame zero are always rebound,
    \item $A_0$ is computed once per contact from incoming velocity and the
    force law,
    \item if a starting overlap has $A_P > A_0$, then $A_0$ is promoted to the
    current penetration area,
    \item one shader invocation owns one source particle and writes only that
    particle's state.
\end{itemize}

Under those assumptions, each source particle needs one contact-state slot for
each particle contact slot in \texttt{ccs}.  The state is source-owned, so a
physical pair can have two mirrored records:

\[
P_i.\texttt{gcs} \rightarrow j
\qquad\text{and}\qquad
P_j.\texttt{gcs} \rightarrow i.
\]

The minimum persistent data is:

\begin{itemize}[leftmargin=*]
    \item whether the slot is active,
    \item target particle index,
    \item phase: compression or rebound,
    \item source velocity at first contact,
    \item target velocity at first contact,
    \item contact normal from the source particle to the target particle,
    \item zero-velocity penetration area $A_0$,
    \item zero-velocity center distance.
\end{itemize}

Everything else can be recomputed from current particle state:

\begin{itemize}[leftmargin=*]
    \item current penetration area $A_P$,
    \item current center distance,
    \item compression fraction $A_P/A_0$,
    \item velocity prediction,
    \item current relative velocity,
    \item reduced mass,
    \item whether the contact has exited.
\end{itemize}

A clear aligned GLSL structure is:

\begin{lstlisting}[language=C]
struct gcontact {
    uvec4 ids;      // x=active, y=target, z=phase, w=unused
    vec4 startVel;  // xy=source first-contact velocity,
                    // zw=target first-contact velocity
    vec4 geom;      // x=normal_x, y=normal_y,
                    // z=A_zero, w=center_distance_zero
};
\end{lstlisting}

The particle structure then stores:

\begin{lstlisting}[language=C]
gcontact gcs[12];
\end{lstlisting}

The twelve slots match the existing twelve \texttt{ccs} particle-contact slots.
Wall state can be added later with four additional slots if wall dynamics are
moved into the same source-owned particle state.

The memory cost is:

\[
16\text{ bytes} + 16\text{ bytes} + 16\text{ bytes}
= 48\text{ bytes per contact}.
\]

For twelve contacts:

\[
12 \times 48 = 576\text{ bytes per particle}.
\]

This memory buys persistent first-contact velocity, phase, and $A_0$.  Those
are the values needed to prevent a stateless GLSL model from locking at the
zero-velocity turnaround point.

\section{Stiffness and Velocity at Overlap}

In the old geometric model, stiffness did not directly change velocity each
frame.  Instead, stiffness determined the penetration depth where the relative
normal velocity should reach zero.

\[
\alpha_0 =
\left(
    \frac{3 \mu v_{\mathrm{in}}^2}{2 q}
\right)^{1/3}
\]

where:

\begin{itemize}[leftmargin=*]
    \item $\alpha_0$ is the zero-velocity penetration depth,
    \item $\mu = m_s m_t / (m_s + m_t)$ is the reduced mass,
    \item $v_{\mathrm{in}}$ is the incoming relative normal speed,
    \item $q$ is the contact stiffness.
\end{itemize}

The zero-velocity center distance is:

\[
d_0 = r_s + r_t - \alpha_0
\]

and the zero-velocity overlap area is:

\[
A_0 = A(d_0)
\]

For the current frame, the compression fraction is:

\[
f_c = \mathrm{clamp}\left(\frac{A}{A_0}, 0, 1\right)
\]

The velocity fraction remaining during compression is:

\[
f_v = \sqrt{1 - f_c}
\]

and the progress toward the zero-velocity point is:

\[
p_c = 1 - f_v
\]

The current compression velocity is then interpolated between the starting
velocity and the zero-velocity-point velocity:

\[
v(A) = v_{\mathrm{start}} +
p_c \left(v_0 - v_{\mathrm{start}}\right)
\]

In short:

\[
\text{stiffness}
\rightarrow
\alpha_0
\rightarrow
A_0
\rightarrow
A/A_0
\rightarrow
v(A)
\]

Higher stiffness produces a smaller $\alpha_0$, so the same overlap removes
more velocity.

\subsection{Chase Collision Problem Area}

The problem area for chase-type collisions is the definition of:

\[
\alpha_0 =
\text{penetration depth where relative normal velocity reaches zero}
\]

In a chase chain, the relative normal velocity of an already-active contact can
change because another contact forms behind one of the particles.  If
$\alpha_0$ is computed once from the original incoming relative velocity and
then treated as fixed for the whole contact episode, the stored zero point can
stop matching the current collision state.

This is the core ambiguity: $\alpha_0$ is both a geometric target and a value
derived from a velocity that may change during chained contacts.



\end{document}
